\documentclass[10pt]{article}
\providecommand{\FIGDIR}{fig}
\newcommand{\DIGESTDATE}{26 September 2026}
\newcommand{\PERIODSTART}{20 September 2026}
\newcommand{\PERIODEND}{26 September 2026}
\input{preamble}

\begin{document}

% ==================================================================== MASTHEAD
\thispagestyle{empty}
\begin{tikzpicture}[remember picture, overlay]
  \fill[Deep] (current page.north west) rectangle
      ($(current page.north east)+(0,-67mm)$);
  \fill[Teal] ($(current page.north west)+(0,-67mm)$) rectangle
      ($(current page.north east)+(0,-68.4mm)$);
  \foreach \x/\y/\r in {22/12/0.9, 41/26/1.5, 63/9/0.7, 88/31/1.1, 112/17/1.9,
                        138/29/0.8, 159/13/1.3, 178/54/1.0, 196/21/1.6,
                        30/44/0.6, 74/57/1.2, 125/48/0.9, 168/6/0.8, 191/59/0.7,
                        52/18/1.0, 100/60/0.6, 148/20/1.1, 66/49/1.4, 205/45/0.8}
    \fill[Sky, opacity=0.30] ($(current page.north west)+(\x mm,-\y mm)$) circle (\r mm);
\end{tikzpicture}

\vspace*{4mm}
{\sffamily\fontsize{7.6}{9}\selectfont\color{Sky}%
 WEEKLY ROLLUP\;\textbullet\;PhD RESEARCH WATCH\par}
\vspace{2.2mm}
{\sffamily\bfseries\fontsize{27}{29}\selectfont\color{white}%
 Particulate Matter\par}
{\sffamily\fontsize{27}{29}\selectfont\color{Sky}%
 Week in Review\par}
\vspace{5.0mm}
{\sffamily\fontsize{8.4}{10}\selectfont\color{white}%
 \mbox{Week ending \DIGESTDATE}\;\;\textcolor{Teal}{\textbar}\;\;\mbox{\PERIODSTART{} to \PERIODEND}\;\;
 \textcolor{Teal}{\textbar}\;\;\mbox{162 records pooled from 7 daily issues}\par}
\vspace{18mm}

% -------------------------------------------------------------- metric strip
\noindent
\begin{minipage}[t]{0.190\linewidth}\metric{162}{RECORDS POOLED\\(162 UNIQUE DOIs)}{Deep}\end{minipage}\hfill
\begin{minipage}[t]{0.190\linewidth}\metric{33\%}{MEASUREMENT SIDE\\(54 RECORDS)}{Teal}\end{minipage}\hfill
\begin{minipage}[t]{0.190\linewidth}\metric{10}{SUBTOPICS\\ACTIVE}{Sky}\end{minipage}\hfill
\begin{minipage}[t]{0.190\linewidth}\metric{31}{POOLED EFFECT\\ESTIMATES (31 WITH A CI)}{Amber}\end{minipage}\hfill
\begin{minipage}[t]{0.190\linewidth}\metric{34}{ABSTRACT-FREE\\(21\% OF WEEK)}{Clay}\end{minipage}

\vspace{6mm}

% -------------------------------------------------------------- week in one box
\section*{\sffamily\bfseries\color{Deep}The week in one paragraph}
\vspace{-1.5mm}{\color{Teal}\rule{\linewidth}{1.1pt}}\vspace{2.5mm}

\begin{keybox}
{\sffamily\bfseries\small\color{Deep}The biggest week in the series by volume, and the worst
since August for missing abstracts: 162 records, 34 of them (21\,\%) carried on metadata alone
--- double last week's rate. The measurement side fell to a series-low third of the week while
\textit{Occupational \& indoor} nearly doubled; the indoor turn named a trend last week is now
the week's dominant movement.}\\[1.4mm]
\footnotesize
\textbf{Scale.} \textbf{162 records} over seven issues on daily totals of \textbf{12}, 18, 23, 27,
\textbf{35}, 24 and 23 --- a 2.9$\times$ range, wider than last week's 1.9$\times$ because Sunday
20 September returned 12. Against recent rollups: 162 after 148, 126, 150, 109 and 112 --- the
highest weekly total recorded. All \textbf{ten} subtopics are populated.

\textbf{The measurement side fell to 33\,\%.} \textit{Exposure assessment \& modelling} eased
38 $\rightarrow$ \textbf{34} and \textit{Sensing} fell 27 $\rightarrow$ \textbf{20}, so the
measurement-side share reads 47, 54, 40, 44, \textbf{33\,\%} across five weeks --- the lowest in the
series. It is not a contraction of the instrumentation literature: \textbf{21 of the 54
measurement-side records are metadata-only} (12 in \textit{Exposure}, 9 in \textit{Sensing}), and
the same Elsevier and ACS deposit pattern that hollowed out those rows also drove the total up. The
share will be re-read when the retry list clears.

\textbf{\textit{Occupational \& indoor} 12 $\rightarrow$ 22.} Last week this watch said the
indoor turn had qualified as a trend and asked whether any of it would move from chambers and
models into occupied buildings with health outcomes. \textbf{Some did.} The 25 September issue
carried a \textbf{948-home} year-round rural heating study (annual mean
\textbf{109\,$\pm$\,207\,$\mu$g\,m$^{-3}$}) and a Xuanwei stove intervention; 24 September a
purifier crossover with mitochondrial DNA methylation outcomes; 26 September a cystic-fibrosis
home-sensor pilot in which \textbf{time above 35\,$\mu$g\,m$^{-3}$}, not the mean, carried the
cytokine signal. Occupied-building work with a biological outcome now exists; none of it is yet
randomised with a clinical endpoint except the \textbf{AIRVAC} protocol entering the registry.

\textbf{Abstract-free doubled, and one reason was an outage.} \textbf{34 of 162 (21\,\%)}
against 15 of 148 (10\,\%). Europe PMC returned HTTP\,503 for the whole 25 September run, removing
both a harvest leg and the second abstract-recovery route, and Elsevier's same-day deposits
without abstracts continued. The deferred-retry pass named as the only remaining fix in the last
two rollups is \textbf{still not built}, and the held list has grown past forty.

\textbf{Fewer estimates, all with intervals.} \textbf{31 pooled estimates}, one per 5.2 records
(series 2.9, 4.4, 2.7, 3.1, \textbf{5.2} --- the sparsest week), but \textbf{31 of 31 carry a CI},
the first complete week. Only the 25 September issue reported none.
\end{keybox}

\clearpage

% -------------------------------------------------------------- trend
\section*{\sffamily\bfseries\color{Deep}Subtopic trend and inflection}
\vspace{-1.5mm}{\color{Teal}\rule{\linewidth}{1.1pt}}\vspace{2.5mm}

\noindent\includegraphics[width=\linewidth]{\FIGDIR/w1_subtopic_trend.png}
\figcap{Per-issue subtopic counts across the week from \texttt{state/metrics.csv}, seven issues.
Daily totals 12, 18, 23, 27, 35, 24 and 23 --- a 2.9$\times$ range. The Sunday low (20 September)
returns after a week's absence; the 35-record Thursday (24 September) carried the CHARLS
constituent and purifier-methylation records and is a complete day, not a double build. Four of
the week's seven issues were built late (the 22:00 rule deferred each by up to a day), none left a
gap.}

\footnotesize
\textbf{Expanding.} \textit{Occupational \& indoor} \textbf{12 $\rightarrow$ 22} is the week's
largest movement (above). \textit{Mechanistic toxicology} \textbf{14 $\rightarrow$ 20} continues a
third week of recovery, led by nitrated-PM inflammation, size-resolved constituent toxicity and
the iron--superoxide results. \textit{Cardiovascular \& metabolic} \textbf{11 $\rightarrow$ 15} and
\textit{Burden, policy \& mitigation} \textbf{13 $\rightarrow$ 17} both rise;
\textit{Reproductive \& developmental} \textbf{5 $\rightarrow$ 8} leaves the bottom of the table
for the first time in this series, carried by the Butajira fetal-biometry cohort, the Chengdu
monitor-versus-satellite preterm analysis, the Chiang Rai placental microscopy and the Shenzhen
Birth Cohort profile.

\vspace{1.2mm}
\textbf{Contracting.} \textit{Sensing} \textbf{27 $\rightarrow$ 20} and \textit{Exposure}
\textbf{38 $\rightarrow$ 34}, both inflated in metadata-only share; \textit{Respiratory \&
allergic} \textbf{13 $\rightarrow$ 10}. \textit{Neuro} (7 $\rightarrow$ 8) and \textit{Other
clinical} (8 $\rightarrow$ 8) hold.

\vspace{1.2mm}
\textbf{The inflection worth naming.} \textbf{Black carbon led the constituent weighting in both
cardiovascular constituent analyses of the week} --- CHARLS incident stroke (per IQR BC \textbf{HR
1.77, 1.54--2.03}, above PM$_{2.5}$ mass at 1.55) on 24 September, and the Kailuan heart-failure
cohort (BC \textbf{40.2--48.7\,\%} of the mixture weight) on 25 September. Two independent Chinese
cohorts, two cardiovascular outcomes, one constituent. For a sensor watch this is the most
actionable signal of the month: BC is the one health-leading constituent that a low-cost node can
measure (filter-based absorption), and mass-only networks cannot see it. The test for next week is
whether a non-Chinese cohort shows the same ordering.

\vspace{2mm}
\noindent\includegraphics[width=\linewidth]{\FIGDIR/f2a_architecture.png}
\figcap{Study architecture over the pooled week. \textbf{\textit{Metadata only} is the largest
single group at 34}, ahead of \textit{Measurement campaign} (29) and \textit{Experimental /
toxicology} (20) --- the first week since August in which the missing-abstract bucket leads.
Primary observational epidemiology (cohort 16, cross-sectional 14, acute 5, ecological 5) totals
\textbf{40}. \textbf{Four} interventional records, against one last week: the AIRVAC and COATED-AIR
trial designs, the Xuanwei stove intervention and a synbiotic RCT.}

\clearpage

% -------------------------------------------------------------- pooled analytics
\section*{\sffamily\bfseries\color{Deep}Pooled analytics}
\vspace{-1.5mm}{\color{Teal}\rule{\linewidth}{1.1pt}}\vspace{3mm}

\noindent\includegraphics[width=\linewidth]{\FIGDIR/f1_subtopics.png}
\figcap{Pooled subtopic distribution, \textbf{all ten bands populated}. \textit{Exposure} 34 and
\textit{Sensing} 20 put the measurement side at \textbf{54 of 162 (33\,\%)}. \textit{Occupational
\& indoor} at 22 is now the largest health-side band, ahead of \textit{Mechanistic toxicology} and
\textit{Sensing} at 20 and \textit{Burden} at 17.}

\vspace{2mm}
\noindent\includegraphics[width=\linewidth]{\FIGDIR/f2b_endpoint.png}
\figcap{Health endpoint over the pooled week, canonicalised. \textbf{``No health endpoint''
pools 78 of 162 (48\,\%)} --- missing abstracts (24) plus exposure, instrument, chemistry and
emission records --- up from 41\,\% last week, the abstract gap showing through. Among human
outcomes, \textbf{respiratory leads at 13}, cardiovascular 10, mortality/longevity 5, then a run of
labels at 4 (cerebrovascular, inflammatory biomarkers, occupational, reproductive, neuro/cognitive,
other clinical).}

\vspace{2mm}
\noindent\includegraphics[width=\linewidth]{\FIGDIR/f3_metric_geography.png}
\figcap{Left --- particle metric. \textbf{PM$_{2.5}$ alone 65 of 162}, unspeciated PM or emissions
61, joint PM$_{2.5}$ + PM$_{10}$ 19. The tail: optical properties 6, \textbf{composition 4},
PM$_{10}$ alone 3, \textbf{size-resolved 2 and ultrafine 1}, bioaerosol 1. Right --- geography.
\textit{Global / multi-region} 66 overstates reach, as every week, pooling reviews, laboratory
studies and unrecoverable metadata-only records. \textbf{China 40} is the largest single-country
bar, then Europe 14, North America 12 and East Asia (ex-China) 8; Sub-Saharan Africa reached
\textbf{5} (Luanda, Durban, Butajira among them).}

\vspace{2mm}
\noindent\includegraphics[width=\linewidth]{\FIGDIR/f6_lifecourse.png}
\figcap{Life-course windows across the pooled week; counts non-exclusive. Working age
\textbf{19}, preconception / in utero \textbf{10} and older adults \textbf{10}, infancy 6.
\textbf{Adolescence 3} --- still the weakest window, carried by APEAL spirometry, a synbiotic
RCT and the AIRVAC protocol (also counted in childhood).}

\clearpage

\noindent\includegraphics[width=\linewidth]{\FIGDIR/f5_forest.png}
\vspace{1mm}
\noindent\includegraphics[width=\linewidth]{\FIGDIR/f5_forest_2.png}
\figcap{All \textbf{31} pooled estimates, ratio and difference scales in separate panels.
\textbf{Not mutually comparable and not meta-analysed}: contrasts include per-IQR, per-SD,
per-10 and per-1\,$\mu$g\,m$^{-3}$, per-doubling, dichotomised and unstated increments, on HR,
sHR, OR, RR and RR-equivalent scales; heterogeneity is structural, so no $I^2$ is reported.
Readable points: \textbf{the largest ratio is the least precise} --- \textbf{OR 4.78
(1.61--14.29)} per doubling of same-day PM$_{2.5}$ for type A aortic dissection, from
\textbf{40 cases}; \textbf{the tightest interval is the smallest effect} --- 0.9977
(0.9972--0.9983) scholar publications per 1\,$\mu$g\,m$^{-3}$ annual PM$_{2.5}$; three CHARLS
analyses (stroke $\times$2, CVD) share participants and are not independent; and the Taiwan COPD
case-control shows the exposure-window dependence directly (1-month OR 1.252, 12-month 1.642,
same cases). The difference panel carries the Butajira fetal-biometry deficits (femur length
\textbf{$-$0.43\,SD} per IQR) and the Bangkok rhinitis VAS slope.}

\vspace{2mm}
\noindent\includegraphics[width=\linewidth]{\FIGDIR/f4_heatmap.png}
\figcap{Subtopic $\times$ architecture, pooled. The densest cell is \textit{Mechanistic
toxicology} $\times$ \textit{Experimental} (\textbf{14}), then \textit{Exposure} $\times$
\textit{Measurement campaign} (13), \textit{Exposure} $\times$ \textit{Metadata only} (12) and
\textit{Sensing} $\times$ \textit{Metadata only} (9) --- the last two are the abstract gap
concentrated on the measurement side. The standing finding survives a ninth week: of the
\textbf{54} measurement-side records, \textbf{4 carry any outcome label and 1 observes a human
outcome} (the PMRI ecological mortality preprint); two are modelled risk indices and one is an
ecotoxicity assay.}

\vspace{4mm}

% -------------------------------------------------------------- venues + trials
\section*{\sffamily\bfseries\color{Deep}Venues, trial watch, and what the pipeline did}
\vspace{-1.5mm}{\color{Teal}\rule{\linewidth}{1.1pt}}\vspace{2.5mm}

\footnotesize
\textbf{Venues.} A three-way tie at \textbf{9} --- \textit{Environ Sci Technol}, \textit{Environ
Pollut}, \textit{Atmos Environ} --- then \textit{Atmos Pollut Res} and \textit{Toxics} at 7,
\textit{Atmos Chem Phys}, \textit{Environ Health (Wash)} and \textit{Front Public Health} at 5.
The leading venue holds 6\,\%. \textit{Atmos Environ} and \textit{Atmos Pollut Res} are the main
sources of metadata-only records (Elsevier same-day deposits). \textbf{Preprints supplied 6}, up
from 2, all labelled.

\vspace{1.5mm}
\textbf{Recurring authors.} Unchanged: surname frequency is a naming artefact, and \textbf{no
repeat-group cluster is claimed}. The one real repeat is a \emph{dataset}, not a group: CHARLS
supplied three separate PM--cardiometabolic papers in the week.

\vspace{1.5mm}
\textbf{Trial watch --- one new trial, one resolved, one update.} \textbf{NCT07196436 (AIRVAC)}
entered on 25 September via its \textit{PLoS One} protocol: four arms, active or sham energy
recovery ventilator $\times$ active or sham portable air cleaner, $>$80 Dallas--Fort Worth asthma
households; \texttt{analyze\_endpoints} shows \textbf{spirometry and ACT as co-primary} in the
registry against ACT alone in the protocol. The \textbf{COATED-AIR} design paper (22 September) was
resolved to \textbf{NCT06541691} ($n$=3,000, Tehran): the air-pollution randomisation's primary is a
composite of non-fatal ischaemic stroke, type I MI, acute limb events or CV death over 30 months,
with \textbf{no registered exposure or adherence endpoint}, so whether the mask-and-alert package
changes anyone's PM dose will not be known; the record still reads \textit{recruiting} six months
after its listed primary completion. \textbf{NCT02944682 (HAPIN)}, the completed multi-country LPG
trial, posted a registry update. \textbf{NCT05423665} reappeared twice and was already held.
\texttt{trials.json} now holds \textbf{16 trials}.

\vspace{1.5mm}
\textbf{What the harvesting legs did.} \textbf{The PubMed connector was load-bearing every day},
adding PMIDs the harvester's two PubMed queries missed and supplying admissions on each issue. The
harvester's \textbf{PubMed sensing query returned 0 on 20--23 September} and 1, 5 and 2 on the last
three days. \textbf{Crossref-by-ISSN} remained the highest-yield leg on weekdays. \textbf{Europe PMC
was down for the 25 September run} (HTTP\,503) and back with 45 records on 26 September.
\textbf{OpenAlex} contributed nothing all week (429 / paywalled date filter). \textbf{arXiv}
returned HTTP\,406 through 25 September and was answering again by the 29 September run (no entry
dated inside the window). \textbf{Consensus} returned the right literature and \textbf{0 new}
records every day, for the fifth consecutive week. \textbf{Scholar Gateway} failed with an identity
error on its first attempt this week. Proofing caught caption and prose miscounts, one unsourced
number and orphaned headings on several days; all were fixed before shipping.

\vspace{2mm}
% -------------------------------------------------------------- gaps
\section*{\sffamily\bfseries\color{Deep}Open gaps}
\vspace{-1.5mm}{\color{Teal}\rule{\linewidth}{1.1pt}}\vspace{2.5mm}

\footnotesize
\begin{itemize}
\item \textbf{The abstract-retry pass, third week running.} 21\,\% of the week shipped without an
abstract, and the most watch-relevant titles are disproportionately among them (a year-long
personal PM/CO$_2$ monitor pilot, Korean station-relocation discontinuities, an AS\&T real-time EC
monitor comparison, a DMA geometry paper, the NAKO mental-health cohort). This is now the single
largest loss of signal in the pipeline.

\item \textbf{Black carbon is leading cardiovascular constituent analyses; the watch has no BC
sensing record this week with a health link.} The two constituent cohorts use modelled BC; none of
the week's absorption-measurement records (Beijing lensing apportionment, Italian 10-year eBC
trends) touches an outcome.

\item \textbf{One of 54 measurement-side records observes a human outcome}, ninth consecutive
week.

\item \textbf{Ultrafine and size-resolved number: 3 of 162}, the same share as last week, despite
a real-time ultrafine size-distribution paper on 22 September.

\item \textbf{Pipeline defects carried forward}: no first-author affiliation field; the
\texttt{Cancer}/\texttt{Oncologic} split in \texttt{ENDPOINT\_CANON}; no \texttt{\textbackslash
DIGESTDATE}-vs-\texttt{--date} build guard; the 2026-07-27 corpus missing \texttt{LIFECOURSE};
\texttt{.git} debris on the delete-restricted mount. \texttt{metrics.csv} continues clean (csv.writer,
500 rows, 0 malformed).
\end{itemize}

\clearpage

% -------------------------------------------------------------- digest
\section*{\sffamily\bfseries\color{Deep}Paper-level digest --- the week by subtopic}
\vspace{-1.5mm}{\color{Teal}\rule{\linewidth}{1.1pt}}\vspace{1mm}

{\fontsize{7.4}{9}\selectfont\color{Slate}Every record that appeared in a daily issue
keeps that issue's critical summary verbatim; nothing here is re-written or re-assessed.
Records surfaced without an abstract are shown as such. Ordered by subtopic, then by the
date of the issue that first carried the record.\par}

\input{digest_body}

% -------------------------------------------------------------- register
\section*{\sffamily\bfseries\color{Deep}Pooled corpus register}
\vspace{-1.5mm}{\color{Teal}\rule{\linewidth}{1.1pt}}\vspace{2.5mm}

\input{table_rows}

\vspace{4mm}
\begin{keybox}
{\sffamily\bfseries\footnotesize\color{Deep}Provenance and method}\\[1.2mm]
{\fontsize{7.2}{8.8}\selectfont
Period \textbf{20--26 September 2026}, Sunday to Saturday, \textbf{seven daily issues},
\textbf{162 records}, \textbf{162 unique DOIs}, no record without a DOI. This rollup is
\textbf{aggregation only}: it re-reads \texttt{state/corpus/*.json}, \texttt{state/metrics.csv} and
\texttt{state/trials.json} over the period and \textbf{queries no literature API}; every record was
gated by \texttt{check\_dois.py} at daily-issue time. The only live queries made for this issue were
the ClinicalTrials.gov reconciliation (COATED-AIR resolution and \texttt{analyze\_endpoints}), which
are state maintenance. Tier mix: \textbf{25 A, 64 B, 73 C}. The W39 weekly was due on the 26
September run; that run and those for 27 and 28 September did not fire, and the rollup was built
on the \textbf{29 September backfill run} after the 26 September daily issue, which it includes.
Effect estimates are transcribed verbatim from the daily issues on heterogeneous contrasts and are
\textbf{not} meta-analysed. Abstracts remain the copyright of their respective publishers.\par}
\end{keybox}

\end{document}
